Los alcanos y alquenos son familias de hidrocarburos que difieren en su estructura de enlace: los primeros poseen enlaces simples carbono--carbono, mientras que los segundos contienen al menos un doble enlace carbono--carbono. Esta diferencia es relevante para interpretar las rutas de transformación reportadas y sus condiciones experimentales. Por otra parte, propiedades macroscópicas como los puntos de fusión y ebullición dependen no solo de la masa molecular, sino también de la estructura y la naturaleza de las interacciones intermoleculares.

Este trabajo estudia metano, etano, propano, butano, isobutano, eteno, propeno, but-1-eno, but-2-eno y pent-1-eno. La fuente primaria de los registros es una recopilación del grupo basada en CAS-SciFinder, organizada en hojas de cálculo~\cite{datoscrudos,datoslimpios,cas}. El interés del análisis no es predecir rendimientos o establecer mecanismos a partir de un número limitado de observaciones, sino caracterizar los registros disponibles, identificar patrones descriptivos y examinar las limitaciones de la información.
